\section{Preliminaries}

\subsection{Notation and Polynomial Encoding}

We work over a finite field $\mathbb{F}$. All vectors, matrices, and tensors are
flattened in row-major order before commitment. For $n \in \mathbb{N}$, we write
$[n] := \{0,1,\ldots,n-1\}$.

All committed lengths are assumed to be powers of two. To make this precise, we
fix an integer $L \ge 1$ such that every committed object is smaller than $2^L$,
and is padded to a length
dividing $2^L$. We also fix a primitive $2^L$-th root of unity $\omega \in
\mathbb{F}$. For every power of two $n$ with $n \mid 2^L$, we define
\[
\omega_n := \omega^{2^L/n}
\qquad\text{and}\qquad
\mathbb{H}_n := \langle \omega_n \rangle.
\]
Thus $\mathbb{H}_n$ is the multiplicative subgroup of size $n$ used as the
evaluation domain for length-$n$ vectors. Its vanishing polynomial is
\[
\nu_n(X) := \prod_{a \in \mathbb{H}_n}(X-a) = X^n-1.
\]
Because all domains are derived from the same root $\omega$, we have
$\omega_{nm}^m=\omega_n$ whenever $n$, $m$, and $nm$ are powers of two dividing
$2^L$. This compatibility is used later when we align objects of different
shapes.

Let $\lambda_n^{(i)}(X)$ denote the $i$-th Lagrange basis polynomial on
$\mathbb{H}_n$.

For a vector $\mathbf{v}\in\mathbb{F}^n$, its polynomial encoding over
$\mathbb{H}_n$ is
\[
v(X) := \sum_{i \in [n]} \mathbf{v}[i] \lambda_n^{(i)}(X),
\]
where $\mathbf{v}[i]$ denotes the $i$-th entry of $\mathbf{v}$.

Throughout the paper, we need to manipulate the shapes
of a high-dimensional tensor. Given $\mathbf{V}\in \mathbb{F}^{a\times b\times c}$,
we write $\mathbf{V}_{\{abc\}}\in \mathbb{F}^{abc}$ for its flatten form,
and $\mathbf{V}_{\{a\times bc\}}\in \mathbb{F}^{a\times bc}$ for its flatten form with the first dimension preserved.
As a sanity check, note that
\[
\mathbf{V}_{\{a\times b\times c\}}[i, j, k] = \mathbf{V}_{\{a\times bc\}}[i, j \cdot c + k] = \mathbf{V}_{\{abc\}}[i \cdot bc + j \cdot c + k].
\]
A similar notation applies to tensors of other dimensions.

\begin{table*}[t]
\centering
\small
\begin{tabular}{lllll}
\toprule
Symbol & Shape & Support & Meaning & Status \\
\midrule
$\mathbf{D}$ & $M\times d$ & all samples & quantized dataset & committed \\
$\mathbf{Y}$ & $M\times K$ & all samples & one-hot labels & committed \\
$\mathbf{S}^{(t)}$ & $M\times K$ & round $t$ & cumulative scores & committed slice \\
$\mathbf{G}^{(t)}$ & $M\times K$ & round $t$ & pseudo-residuals & committed slice \\
$\boldsymbol{\ell}^{(t,k)}$ & $M$ & tree $(t,k)$ & leaf assignment & committed slice \\
$\mathbf{C}^{(t,k)},\mathbf{H}^{(t,k)}$ & $N_{\mathrm{tot}}\times d\times B$ & tree $(t,k)$ & count / sum histograms & committed slice \\
$\Gamma^{(t,k)}$ & $N_{\mathrm{int}}\times d\times B$ & tree $(t,k)$ & candidate split scores & committed slice \\
$\widetilde{\mathbf{W}}^{(t)}$ & $N_{\mathrm{tot}}\times K$ & round $t$ & padded leaf values & committed slice \\
$F(X),\widehat{u}(X)$ & varies & derived & batched virtual polynomials & virtual \\
\bottomrule
\end{tabular}
\caption{Main proof-facing objects used later in Sections~3--5. All tensors are
flattened in row-major order before commitment.}
\label{tab:object-glossary}
\end{table*}

\subsection{KZG Polynomial Commitments}

We use a KZG-based polynomial commitment scheme (PCS) for univariate polynomials.
The PCS consists of the following algorithms:

\begin{itemize}
    \item $\Setup(1^\lambda, d_{\max}) \rightarrow \pp$:
    Generates public parameters supporting commitments to polynomials of degree at most $d_{\max}$.

    \item $\Commit(\pp, f) \rightarrow c$:
    Commits to a polynomial $f$ and outputs a commitment $c$.

    \item $\Open(\pp, f, z) \rightarrow (v, \pi)$:
    Given a polynomial $f$ and an evaluation point $z$, outputs the evaluation
    $v = f(z)$ together with an opening proof $\pi$.

    \item $\Verify(\pp, c, z, v, \pi) \rightarrow \{0,1\}$:
    Verifies that $v$ is the correct evaluation of the polynomial committed in $c$
    at point $z$.
\end{itemize}

\paragraph{Batched openings.}
In addition to single-point openings, we use a batched opening interface that
allows the prover to open multiple evaluation claims together.

\begin{itemize}
    \item $\BatchOpen(\pp, \{f_i\}_{i \in [m]}, \{Q_i\}_{i \in [m]})
    \rightarrow (\{v_{i,z}\}_{i \in [m], z \in Q_i}, \Pi)$:
    For each committed polynomial $f_i$ and each query point set $Q_i$,
    outputs all claimed evaluations $v_{i,z} = f_i(z)$ together with a single
    batched proof $\Pi$.

    \item $\BatchVerify(\pp, \{c_i\}_{i \in [m]}, \{Q_i\}_{i \in [m]},
    \{v_{i,z}\}_{i \in [m], z \in Q_i}, \Pi) \rightarrow \{0,1\}$:
    Verifies that all claimed evaluations are consistent with the corresponding
    commitments.
\end{itemize}

\paragraph{Properties.}
We rely on the following standard properties. \emph{Completeness}: honestly generated openings
and batched openings always verify.  \emph{Binding}: after a
commitment is published, the prover cannot later open it to a different value at
the same point. 

\subsection{Gradient Boosted Decision Trees in Plaintext}

Given a commitment to a forest, the task of proof of training is to show that
the forest arises from a valid execution of the GBDT training procedure.

Let $\mathbf{D}\in[B]^{M\times d}$ be the quantized training dataset, where
$M$ is the number of samples, $d$ is the number of features, and each feature
value lies in one of $B$ bins. Let $\mathbf{Y}\in\{0,1\}^{M\times K}$ be the
one-hot label matrix for a $K$-class classification task. We treat the number
of boosting rounds $T$,
the number of classes $K$,
and the size of dataset $M$
as public parameters of the statement being proved.

\paragraph{Plaintext training.}
GBDT trains an additive model over $T$ rounds. In each round $t\in[T]$, the
learner maintains a score matrix $\mathbf{S}^{(t)}\in\mathbb{F}^{M\times K}$,
where $\mathbf{S}^{(t)}[r,k]$ denotes the current score of sample $r$ for class
$k$, before adding the trees of round $t$. We initialize
$\mathbf{S}^{(0)}=\mathbf{0}$.

The current class probabilities are obtained by applying softmax row-wise:
\begin{equation}
\boldsymbol{\pi}^{(t)}[r,k]
:=
\frac{\exp(\mathbf{S}^{(t)}[r,k])}
{\sum_{k'\in[K]} \exp(\mathbf{S}^{(t)}[r,k'])}.
\end{equation}
Let $\mathbf{Y}\in\{0,1\}^{M\times K}$ be the one-hot label matrix. The
multinomial logistic loss at round $t$ is
\begin{equation}
\mathcal{L}^{(t)}
:=
-\sum_{r\in[M]}\sum_{k\in[K]}
\mathbf{Y}[r,k]\log \boldsymbol{\pi}^{(t)}[r,k].
\end{equation}
Its partial derivative with respect to the score entry
$\mathbf{S}^{(t)}[r,k]$ is
\begin{equation}
\frac{\partial \mathcal{L}^{(t)}}{\partial \mathbf{S}^{(t)}[r,k]}
=
\boldsymbol{\pi}^{(t)}[r,k]-\mathbf{Y}[r,k].
\end{equation}
Therefore, the negative gradient is
\begin{equation}\label{eq:plain-pseudoresidual}
\mathbf{G}^{(t)}[r,k]
:=
-\frac{\partial \mathcal{L}^{(t)}}{\partial \mathbf{S}^{(t)}[r,k]}
=
\mathbf{Y}[r,k]-\boldsymbol{\pi}^{(t)}[r,k].
\end{equation}
We call $\mathbf{G}^{(t)}$ the pseudo-residual matrix. Similar to the gradient descent
approach, we want to move the score matrix towards the negative gradient direction, and therefore reduce the loss.
Thus in round $t$, the
learner trains one regression tree for each class $k\in[K]$ to fit the target
vector $\mathbf{G}^{(t)}[\cdot,k]\in\mathbb{F}^{M}$.\footnote{The above describes the training process for typical gradient-boosting decision trees. 
It is also possible to adapt our constructions to other variants,
e.g. XGBoost~\cite{xgboost16},
by modifying the training target and tree-splitting policy accordingly.
See \textcolor{red}{Appendix}.}

\paragraph{Regression tree.}
Fix a round $t$ and a class $k$. The learner constructs a regression tree
$f^{(t,k)}:[B]^d\to\mathbb{F}$, which is a binary decision tree whose leaves
carry scalar prediction values. Each internal node $u$ is labeled by a feature
index $j_u\in[d]$ and a threshold bin $b_u\in[B]$. Given a sample
$\mathbf{x}\in[B]^d$, routing at node $u$ is determined by the predicate
$\mathbf{x}[j_u] < b_u$: if it holds, the sample is sent to the left child of
$u$; otherwise it is sent to the right child. The output
$f^{(t,k)}(\mathbf{x})$ is the value stored at the leaf reached by this routing
procedure.

Similarly, during training, each node $u$ is associated with the set
$I_u\subseteq[M]$ of samples currently reaching $u$. The root contains all
samples, i.e., $I_{\mathrm{root}}=[M]$, and each split partitions the parent
sample set into the left and right child sample sets.

\paragraph{How a node is split.}
For the class-$k$ tree in round $t$, let
\[
g_r := \mathbf{G}^{(t)}[r,k]
\]
denote the regression target attached to sample $r$. Consider a node $u$ with
current sample set $I_u$. For a candidate split specified by a feature
$j\in[d]$ and a threshold bin $b\in[B]$, define
\begin{align}
L_{u,j,b} &:= \{\, r\in I_u : \mathbf{D}[r,j] < b \,\},\\
R_{u,j,b} &:= \{\, r\in I_u : \mathbf{D}[r,j] \ge b \,\}.
\end{align}
The quality of this split is measured by the residual sum of squares (RSS):
\begin{equation}\label{eq:rss-plain}
\operatorname{RSS}(u,j,b)
=
\sum_{r\in L_{u,j,b}} (g_r-\bar{g}_L)^2
+
\sum_{r\in R_{u,j,b}} (g_r-\bar{g}_R)^2,
\end{equation}
where
\[
\bar{g}_L := \frac{1}{|L_{u,j,b}|}\sum_{r\in L_{u,j,b}} g_r,
\qquad
\bar{g}_R := \frac{1}{|R_{u,j,b}|}\sum_{r\in R_{u,j,b}} g_r
\]
are the empirical means on the two sides. Intuitively, the split is good if the
targets assigned to each child are well concentrated around a single mean.

Among all valid candidates, i.e. those for which both $L_{u,j,b}\neq\emptyset$ and $R_{u,j,b}\neq\emptyset$, the trainer selects a
feature-threshold pair minimizing \eqref{eq:rss-plain}.
If valid split is not possible, the node is not split and becomes a leaf.

Once the tree structure is fixed, each leaf $\ell$ is assigned the mean target
value of the samples that reach it. Writing $I_\ell\subseteq[M]$ for the sample
set of leaf $\ell$, its output value is
\begin{equation}\label{eq:plain-leaf-value}
w_{\ell}^{(t,k)}
:=
\frac{1}{|I_\ell|}\sum_{r\in I_\ell} \mathbf{G}^{(t)}[r,k].
\end{equation}
Therefore, for every training sample $r$, the tree prediction
$f^{(t,k)}(\mathbf{D}[r,\cdot])$ is exactly the leaf value associated with the
leaf reached by sample $r$.

\paragraph{Round update.}
After training all $K$ regression trees
$\{f^{(t,k)}\}_{k\in[K]}$ in round $t$, the learner updates the score matrix by
adding the tree outputs with learning rate $\mu>0$:
\begin{equation}\label{eq:plain-score-update}
\mathbf{S}^{(t+1)}[r,k]
=
\mathbf{S}^{(t)}[r,k] + \mu\, f^{(t,k)}(\mathbf{D}[r,\cdot]).
\end{equation}
This completes round $t$. The next round recomputes
$\boldsymbol{\pi}^{(t+1)}$, forms the new pseudo-residual matrix
$\mathbf{G}^{(t+1)}$ using \eqref{eq:plain-pseudoresidual}, and trains a fresh
collection of $K$ regression trees against those updated targets. After $T$
rounds, the forest consists of the $TK$ trees trained all rounds.

At inference time, for a fresh sample $\mathbf{x}\in[B]^d$, the model
accumulates the per-class scores
\[
\sum_{t=0}^{T-1} f^{(t,k)}(\mathbf{x})
\qquad\text{for each }k\in[K],
\]
and predicts from the resulting $K$-dimensional score vector, typically by
applying softmax.

\paragraph{Implication for the proof.}
The proof-critical step is the correctness of the regression-tree training at
each node, namely, that the committed split metadata and leaf values are
consistent with the pseudo-residual targets of that round and class. In
standard plaintext implementations, the optimizer often evaluates split
candidates by sorting samples according to each feature and scanning candidate
thresholds in that sorted order. In quantized systems, the same optimization is
commonly expressed through per-bin histograms. Our proof follows the latter
view: instead of proving a sorting procedure in zero knowledge, we prove the
correctness of the histogram statistics from which the trainer derives the same
split scores.

\subsection{Quantization}

We follow the fixed-point quantization approach used in prior ZKP for ML works~\cite{sparrow24,zkdt20,zkboost26}. 
Given a real number $x$, we encode it as a signed integer
\[
q_x = \lfloor x \cdot 2^Q \rfloor,
\]
where $Q>0$ is a public quantization parameter. The corresponding fixed-point value is then interpreted as $x \approx q_x / 2^Q$.

Under this representation, addition is scale-preserving: if
$x_1 \approx q_{x_1}/2^Q$ and $x_2 \approx q_{x_2}/2^Q$, then
\[
x_1 + x_2 \approx \frac{q_{x_1}+q_{x_2}}{2^Q}.
\]
Multiplication, however, increases the scale and therefore requires an explicit
rescaling step. In particular, if
$x \approx q_x/2^Q$ and $y \approx q_y/2^Q$, then their product is represented by
\[
q_z = \left\lfloor \frac{q_x q_y}{2^Q} \right\rfloor,
\]
so that
\[
xy \approx \frac{q_z}{2^Q}.
\]
\subsection{Statement and Proof Overview}

Our goal is to prove correct training, not merely well-formedness of the
committed forest. Concretely, given commitments to a training dataset and to a
forest, the prover must show that the committed forest is exactly the result of
running the public quantized GBDT training algorithm on the committed dataset.

The public statement consists of the training parameters
$(M,d,B,K,T)$, the learning rate $\mu$, the public tree-shape policy, the
quantization parameters, the public tables used by imported subprotocols, the
KZG setup parameters, and the commitments to the dataset and the forest. The
witness consists of the values needed to certify one valid training trace,
including the labels, the round-wise boosting states, the histogram
statistics, the split metadata, the leaf values, and the auxiliary witnesses
required by the imported gadgets.

Accordingly, the verifier checks the following claim: there exists a witness
such that, starting from the committed dataset, the public quantized training
procedure produces exactly the committed forest. Thus, an accepting proof
certifies not only that the committed forest has valid syntax, but that it is
consistent with one end-to-end execution of the trainer semantics fixed by this
paper.

At a high level, the proof of training must establish four facts:
\begin{enumerate}
    \item the committed forest derives the claimed inference over the dataset;
    \item each boosting round correctly derives the pseudo-residuals and updates
    the score state;
    \item in each boosting round, each tree is trained according to the public training algorithm.
\end{enumerate}

The rest of the protocol decomposes this global claim into local checks for
forest semantics, routing, histogram construction, split-score correctness,
split optimality, and leaf-value correctness, and then batches these checks
into a small number of polynomial openings.

\begin{theorem}[Informal security statement]
\label{thm:main-informal}
Assume the KZG commitment scheme is binding, and the imported subprotocols satisfy
their claimed completeness and soundness properties. Then the protocol satisfies:
\begin{enumerate}
    \item \textbf{Completeness.} If the committed forest is obtained by honestly
    running the public quantized trainer on the committed dataset, then the
    prover can produce an accepting proof.

    \item \textbf{Soundness.} If the verifier accepts, then, except with the
    accumulated soundness error of the imported gadgets and the batching
    checks, the committed forest is consistent with a valid execution of that
    same trainer on the committed dataset.

    \item \textbf{(Optional) Witness privacy.} If the hiding variant of the commitment
    scheme is used and the imported subprotocols are zero knowledge, then the
    proof reveals nothing beyond the public statement and the committed
    outputs.
\end{enumerate}
We claim soundness, not a separate argument-of-knowledge statement.
\end{theorem}
